\documentclass[12pt]{article}
\usepackage[T1]{fontenc}
\usepackage[utf8]{inputenc}
\usepackage{lmodern}
\usepackage{textcomp}
\usepackage{enumitem}
\usepackage{geometry}
\geometry{margin=1in}
\begin{document}
\begin{center}
Answer Key - Economics - Graduate Level Assessment 13 \
\small (Answers are ordered exactly as in the question paper.)
\end{center}

\section*{Multiple Choice}
\begin{enumerate}[label=\arabic*.]
\item \textbf{Answer:} D
\\ \textit{Options:} A) Excess profits due to lack of competition.; B) Firms produce identical products.; C) High degree of interdependence among firms.; D) Zero economic profits.; E) Large number of buyers and one seller.
\\ \textit{Note:} In the long run, firms in monopolistic competition will enter and exit the market until economic profits are driven to zero due to the presence of free entry and exit, leading to a situation where firms earn just enough to cover their average costs.
\item \textbf{Answer:} A
\\ \textit{Options:} A) Business expectations can affect the economy by affecting the level of investment.; B) Business expectations mold the economy through fluctuations in commodity prices.; C) Business expectations have no impact on the economy.; D) Business expectations steer the economy by directly setting interest rates.; E) Business expectations influence the economy through government policy changes.
\\ \textit{Note:} Business expectations influence the economy because optimistic projections can lead to increased investment in capital and innovation, thereby stimulating economic growth, while pessimistic expectations may result in reduced spending and slower economic activity.
\item \textbf{Answer:} A
\\ \textit{Options:} A) A decrease in the price of orange juice; B) A decrease in the price of Tang; C) An increase in the price of grapefruit juice; D) A decrease in the income of juice drinkers; E) A bumper crop of oranges in Florida
\\ \textit{Note:} A decrease in the price of orange juice typically increases the quantity demanded, which can lead to a rightward shift in the demand curve, reflecting a higher overall demand at that lower price point due to the law of demand.
\item \textbf{Answer:} A
\\ \textit{Options:} A) the real quantity of money increases; B) imports become relatively more expensive; C) the value of cash increases; D) exports become relatively cheaper; E) inflation rate decreases
\\ \textit{Note:} The correct answer is based on the real balance effect, which states that when the price level increases, the purchasing power of money decreases, leading consumers to reduce their spending, thus resulting in a downward sloping aggregate demand curve.
\item \textbf{Answer:} B
\\ \textit{Options:} A) Investment and Consumption; B) Transactions Demand and Asset Demand; C) Income and Expenditure; D) Savings and Interest Rates; E) Credit and Debit
\\ \textit{Note:} Marshall and Fisher identified two primary components of the demand for money: transactions demand, which pertains to the need for money to facilitate everyday transactions, and asset demand, which relates to the desire to hold money as a store of value rather than investing it.
\end{enumerate}

\section*{True / False}
\begin{enumerate}[label=\arabic*.]
\setcounter{enumi}{5}
\item \textbf{Answer:} False
\\ \textit{Options:} A) True; B) False
\\ \textit{Note:} A price ceiling leads to a shortage, not a surplus, because it sets a maximum price that is below the market equilibrium, resulting in higher demand and lower supply.
\item \textbf{Answer:} False
\\ \textit{Options:} A) True; B) False
\\ \textit{Note:} The statement is false because an increase in the price of polyester pants would likely lead to a decrease in the quantity supplied of polyester pants, potentially increasing the demand for denim jeans as a substitute, which would actually raise the price and quantity of denim jeans rather than decrease them.
\item \textbf{Answer:} True
\\ \textit{Options:} A) True; B) False
\\ \textit{Note:} Investment demand tends to increase during a stock market downturn because investors seek safer, more stable investment opportunities, often turning to assets like bonds or real estate, which are perceived as less risky compared to volatile stocks.
\item \textbf{Answer:} False
\\ \textit{Options:} A) True; B) False
\\ \textit{Note:} Structural unemployment is primarily caused by long-term changes in the economy, such as technological advancements and industry shifts, rather than short-term fluctuations in the business cycle.
\item \textbf{Answer:} False
\\ \textit{Options:} A) True; B) False
\\ \textit{Note:} The circular flow model of pure capitalism primarily illustrates transactions between households and firms, excluding direct consumer-to-consumer transactions which are not represented in this framework.
\end{enumerate}

\section*{Long Form Response}
\begin{enumerate}[label=\arabic*.]
\setcounter{enumi}{10}
\item \textbf{Answer:} Anticipated inflation can benefit certain institutions and individuals, notably bank owners and the United States Treasury, by altering the dynamics of borrowing and lending. Bank owners often gain from inflation because they can repay fixed-rate loans with money that is worth less in real terms, thereby increasing their profit margins. Similarly, the U.S. Treasury benefits as inflation can reduce the real burden of national debt; as prices rise, the value of tax revenues increases, allowing the government to service its debt more easily. Overall, these mechanisms highlight how anticipated inflation can create advantageous conditions for specific actors in the economy, enabling them to navigate the challenges posed by rising prices effectively.
\\ \textit{Note:} correct because it illustrates how anticipated inflation allows bank owners to benefit from repaying fixed-rate loans with devalued currency, thereby enhancing profitability, while also demonstrating that the U.S. Treasury gains from a reduced real debt burden as rising prices increase the nominal value of tax revenues.
\item \textbf{Answer:} Logit and probit models differ primarily in their methods for transforming probabilities, which affects their application in economic analysis. The logit model employs the logistic function, which maps a linear combination of independent variables to a probability range between zero and one using the logistic curve. In contrast, the probit model utilizes the cumulative distribution function of the standard normal distribution for the same purpose. These transformations lead to distinct interpretations of the coefficients and influence how the models handle extreme probabilities. Consequently, the choice between logit and probit can impact the estimated probabilities and their interpretation in economic contexts, particularly when analyzing binary outcomes or decision-making scenarios.
\\ \textit{Note:} correct because it accurately identifies the logit model's use of the logistic function and the probit model's reliance on the standard normal distribution's cumulative distribution function, highlighting how these differing transformations of probabilities influence their applicability in economic analysis.
\end{enumerate}

\vfill
\noindent\textit{End of Answer Key}
\end{document}